\documentclass[11pt]{article}
\usepackage[utf8]{inputenc}
\usepackage[a4paper,margin=2.5cm]{geometry}
\usepackage{booktabs,array,longtable,xcolor,hyperref}
\definecolor{Navy}{HTML}{0B1F33}
\definecolor{Blue}{HTML}{16697A}
\hypersetup{colorlinks=true,urlcolor=Blue}
\setlength{\parindent}{0pt}
\setlength{\parskip}{0.6em}

\begin{document}

{\LARGE\bfseries\color{Navy}Guion de presentación del proyecto final de Redes}\par
{\large Fases 1 y 2 · Grupo 1 · Sección 10}

\section*{Uso del guion}

El guion sigue las doce diapositivas de \texttt{presentacion.html}. Cada integrante debe explicar
las decisiones que implementó y reconocer la diferencia entre validación local y evidencia real en
AWS. La duración sugerida es de 12 a 15 minutos, más preguntas.

\begin{longtable}{@{}p{1.1cm}p{3.6cm}p{2.8cm}p{7.1cm}@{}}
\toprule
\textbf{N.} & \textbf{Tema} & \textbf{Presenta} & \textbf{Puntos de explicación} \\
\midrule
1 & Portada & Pablo & Identificar el curso, el equipo y el alcance acumulativo de las dos fases. \\
2 & Objetivo & Pablo & Explicar la necesidad de segmentar usuarios, servicios, administración y exposición pública mediante infraestructura reproducible. \\
3 & Fase 1 & Andrés Rafael & Resumir los 20 recursos desplegados y las dos pruebas ICMP conservadas como evidencia de la base. \\
4 & Arquitectura Fase 2 & Pablo & Recorrer la DMZ, Data Center, clientes y dependencias desplegados durante la ventana temporal. \\
5 & Segmentación & Pablo & Justificar CIDR, capacidad, reserva y direcciones fijas de los servicios. \\
6 & Seguridad & Genser & Diferenciar Security Groups, NACL y nftables. Explicar retornos efímeros e ICMP. \\
7 & DNS e intranet & Andrés Rafael & Explicar BIND9, el nombre interno, redes autorizadas y el acceso HTTP a Recursos Humanos. \\
8 & Web pública & Roberto & Explicar DMZ, Elastic IP temporal, Nginx, hardening y HTTPS opcional. \\
9 & Visitas y TI & Genser & Describir el aislamiento de Visitas y la política asimétrica de TI. Mencionar la instalación privada de nftables y la autoprueba DNS. \\
10 & Validación & Pablo & Presentar los validadores locales y el resultado AWS: 36 correctas, 3 omitidas y 0 contradictorias. \\
11 & Ventana AWS & Andrés Rafael & Explicar la ejecución, identidad, permisos, evidencias, inventario y retiro; Pablo revisó e integró la documentación. \\
12 & Equipo y cierre & Pablo & Presentar al equipo con sus carnés confirmados y recordar que HTTPS permanece como punto extra opcional. \\
\bottomrule
\end{longtable}

\section*{Preguntas previsibles}

\textbf{¿Por qué no se utilizó NAT Gateway?} Las subredes privadas no requieren Internet general.
Los paquetes de Amazon Linux llegan mediante un Gateway Endpoint de S3, que limita el acceso al
repositorio y evita el costo horario de NAT.

\textbf{¿Por qué Security Groups y NACL?} Los Security Groups aplican reglas stateful por interfaz.
Las NACL agregan control stateless por subred y obligan a declarar el tráfico de retorno. nftables
aporta una tercera capa en el sistema operativo.

\textbf{¿Cómo se demuestra el aislamiento?} La suite ejecuta pruebas positivas y negativas desde
el segmento correspondiente. Un resultado solo cuenta como evidencia cuando registra fecha,
commit, origen, destino, comando, esperado, obtenido y veredicto.

\textbf{¿Está terminada la Fase 2?} La implementación base, el despliegue temporal, las pruebas
aplicables, la captura dedicada de la web y los datos de los integrantes están completos. HTTPS
permanece como punto extra opcional.

\textbf{¿Quién tomó las pruebas?} Andrés Rafael Chivalán ejecutó la ventana y recopiló las
capturas; Pablo Daniel Barillas Moreno revisó e integró la documentación. Genser conserva la
autoría de clientes, sondas y automatización; Roberto conserva la autoría de la web pública.

\section*{Datos que deben confirmarse antes de entregar}

\begin{itemize}
  \item Decisión sobre el punto extra de HTTPS.
  \item Aceptación docente de Security Groups más NACL como firewall y ACL del ejercicio.
\end{itemize}

\end{document}
